% Job posting: https://picovoice.ai/careers/deep-learning-researcher/
\documentclass[11pt]{article}

% --------- PACKAGES ----------
\usepackage[left=0.8in,top=0.7in,right=0.8in,bottom=0.7in]{geometry}
\usepackage{enumitem}
\usepackage[hidelinks]{hyperref}
\usepackage{setspace}
\usepackage{parskip}
\usepackage[en-GB]{datetime2}
\usepackage[T1]{fontenc}
\usepackage{newtxtext,newtxmath}   % Times-like
% pdfTeX-only: improves ATS text extraction. Guarded so XeTeX/Tectonic skip it.
\ifdefined\pdfgentounicode
  \input{glyphtounicode}
  \pdfgentounicode=1
\fi

\setlength{\parindent}{0pt}
\setstretch{1.05}
\pagenumbering{gobble}

% --------- FIXED IDENTITY (KG §1 — do not edit) ----------
\newcommand{\clName}{Nishal Stanislaus Silva}
\newcommand{\clEmail}{nishal.silva@hotmail.com}
\newcommand{\clPhone}{+1 613 200 2030}
\newcommand{\clCity}{Toronto, ON}
\newcommand{\clSite}{https://nishal.xyz}
\newcommand{\clLinkedIn}{https://www.linkedin.com/in/nishal-silva}
\newcommand{\clStatus}{Canadian Permanent Resident, Eligible to work in Canada without sponsorship}

% --------- PER-APPLICATION VARIABLES ----------
\newcommand{\clTagline}{Deep Learning Researcher | Real-Time Inference, Time-Series, Pattern Detection | Applied ML}
\newcommand{\clRole}{Deep Learning Researcher}
\newcommand{\clCompany}{Picovoice}
\newcommand{\clTeam}{the research team}
\newcommand{\clAddressee}{Dear Hiring Manager,}
\newcommand{\clReqID}{}

\begin{document}
\pagestyle{empty}

\begin{center}
    {\LARGE \scshape \textbf{\clName}}{\large \textit{, PhD}}
    \\[1pt]
    {\small \clTagline}
    \\[1pt]
    {\footnotesize\textit{\clStatus}}
    \\[1pt]
    {\small
    \href{mailto:\clEmail}{\clEmail}
    \,\,|\,\, \clPhone
    \,\,|\,\, \clCity
    \,\,|\,\, \href{\clSite}{nishal.xyz}
    \,\,|\,\, \href{\clLinkedIn}{linkedin.com/in/nishal-silva}}
\end{center}

\rule{\linewidth}{0.4pt}\vspace{0.6em}

\today

\textbf{Re: \clRole{} -- \clCompany{}}

\clAddressee

I am applying for the Deep Learning Researcher role at Picovoice. I am a Machine Learning Engineer with a PhD (University of Trento, 2025) and 10+ years across industry and academia, and my specialty is deep learning that runs in real time on small devices. Picovoice builds on-device voice AI, and that is the constraint I have worked under for most of my research: models that must be accurate and fast on limited compute.

Three parts of your posting match my work directly. For deep learning on constrained hardware, my published audio pattern detection system (Journal of the Audio Engineering Society, 2026) is a recurrent network that runs in 14 ms on a Raspberry Pi 4 at F1 0.76, 74x faster than the DTW baseline, and I owned it from data acquisition to deployment. For probability and algorithms, I compared probabilistic and deterministic pattern detection methods (FRUCT 2020) and wrote Nebula, an open-source C++17 library of audio features (MFCC, STFT, pitch) with per-feature latency benchmarks on ARM. For research method, I design benchmark suites with baselines and ablations, and I released four open datasets on Zenodo.

I should be direct about the gaps. I have not written CUDA kernels, and my C work is lighter than my C++ and Python. I use PyTorch alongside TensorFlow and Keras, which carried most of my published work, and I have not done formal graph theory work beyond the state machines in my gesture and pattern detection systems. I have shipped real-time C++ and Python inference on embedded hardware, and I expect to pick up GPU-level work quickly. What draws me to Picovoice is that the research is measured by what runs on the device.

I am a Canadian permanent resident based in \clCity{}, eligible to work without sponsorship, and available immediately. I am open to relocating to Vancouver. My portfolio, publications and open datasets are available at \href{\clSite}{nishal.xyz}.

\vspace{12pt}
Sincerely,\\[2pt]
\textbf{\clName}, PhD

\end{document}
